% ===================================================================
%  TABELO N° 10 — La maro. La portuo.
%  Feuillets 70 a 75 du fac-simile = folios 68 a 73.
%  Correspondance : folio imprime = numero de feuillet - 2.
%  Feuillet 76 (folio 74) resta blanka : ol esas la recto vacua avan
%  la TABELO N° 11 (Triesma Serio), qua komencas ye feuillet 77.
% ===================================================================

% --- Feuillet 70 (folio 68) -----------------------------------------
\begin{VUpage}[70]{68}
%%K t10-apar-1 apar
\VUsaut{4.0mm}
\VUtitre{15.0pt}{60}{TABELO N\textsuperscript{o} 10}
\VUsaut{3.0mm}
%  La ligne qui suit est le TITRE du tableau, sous son numero, comme aux
%  tableaux 1 a 6 : le volume francais la garde d'ailleurs dans son bloc
%  d'ouverture, « La Mer. --- Le Port. ». Une cle l'en separait, et la
%  page de lecture la posait donc un rang plus bas, seule, en face d'une
%  case vide -- et en petites capitales de section, non en titre.
\VUpk{11.4pt}{40}{La Maro. --- La Portuo.}
\VUsaut{3.0mm}

%%K t10-01-1 p
\VUblancAlinea
1. --- Somere, dum ke uli serchas kalmeso\nl
e koldeteso sur monto, altri preferas adirar la\nl
litoro. Tale ni agis lastayare, nam ni tre pri\cc
zas \VUgras{Oceano} \textsuperscript{(1)}.

%%K t10-02-1 p
\VUblancAlinea
2. --- Pri me, me sentas granda plezuro kon\cc
templante ta imensa vasteso di aquo qua su\nl
extensas til la \VUgras{horizonto} \textsuperscript{(2)}, e vidante departar\nl
por longa \VUgras{pasajo} l'enorma \VUgras{vapor-navi} \textsuperscript{(3)} sur\nl
qui embarkas la voyajeri iranta aden Ame\cc
rika.

%%K t10-03-1 p
\VUblancAlinea
3. --- Pro to, mea patrulo, qua konocas mea\nl
inklineso, sempre selektas, por pasar la va\cc
kanci, ul plajo tre kalma, situita proxim un\nl
de nia granda \VUgras{milital portui}.

%%K t10-03-2 p
\VUblancAlinea
Dum nia lasta sejorno, l'\VUgras{eskadro} venis an\cc
kragar dop l'enorma \VUgras{digo} \textsuperscript{(4)} qua klozas e pro\cc
tektas la \VUgras{milital portuo} \textsuperscript{(5)}, e me sucesis segun\nl
mea deziro admirar la diversa \VUgras{navi militala},\nl
la mikra \VUgras{submara navo} \textsuperscript{(10)} qua plunjas e des\cc
aparas sub l'aqui ed anke l'enorma \VUgras{kuras}\cc
\VUgras{navo} \textsuperscript{(9)} qua, til nun timis preske nur la ataki\nl
dal \VUgras{torped-navo} \textsuperscript{(8)}.

%%K t10-tit-2 sub
\VUsaut{3.0mm}
\VUpk{10.2pt}{40}{La mikra navi.}
\VUsaut{3.0mm}

%%K t10-04-1 p
\VUblancAlinea
4. --- Olca ja savis tre \VUgras{habile} trovar la febla\nl
punto di la dika metal-kuraso por stekar ibe\nl
la danjeroza \VUgras{torpedo}, di qua l'explozo efekti\cc
\parplein
\end{VUpage}

% --- Feuillet 71 (folio 69) -----------------------------------------
\begin{VUpage}[71]{69}
%%K t10-04-1 p suite
\VUcontinue
gas la periso dil navo, ma tamen ol esis min\nl
timenda kam la submara navo, nam la kontre\cc
torpedieri facile persequas olu.

%%K t10-05-1 p
\VUblancAlinea
5. --- Me anke povis vidar la rapida \VUgras{kroz}\cc
\VUgras{navi} \textsuperscript{(6)} kurasizita, preske tam ne-vundebla kam\nl
la vera kuras-navo, plura \VUgras{transporto-navi} \textsuperscript{(12)}\nl
tri o quar \VUgras{kanon-bateli} \textsuperscript{(7)} e fine un \VUgras{fregato} \textsuperscript{(11)}.\nl
\VUgras{La spektaklo di ta navaro, kande ol manovras}\nl
\VUgras{por desankragar esas maxim interesanta.}

%%K t10-06-1 p
\VUblancAlinea
6. --- Quanta difero di konstrukteso inter\nl
ta granda stal-amasi e la eleganta \VUgras{trimasti} o la\nl
gracioza \VUgras{seglo-navi} \textsuperscript{(13)} ankore nun uzata en\nl
la komercala navaro quan me vidis enirar la\nl
portuo, per tota segli, kande me promenis sur\nl
la \VUgras{warfo} \textsuperscript{(14)}. Advere olci perdis multa de sua\nl
avantaji \VUgras{kande} la \VUgras{vento} esis kontrea. Li mustis\nl
des-hisar omna segli por ri-enirar la baseno e\nl
\VUgras{remorkero} \textsuperscript{(16)} esis necesa por querar li ed ad\cc
duktar li, quale simpla \VUgras{kargo-bateli} \textsuperscript{(17)}, ad la\nl
kayo.

%%K t10-tit-3 sub
\VUsaut{3.0mm}
\VUpk{10.2pt}{40}{La komercala portuo.}
\VUsaut{3.0mm}

%%K t10-07-1 p
\VUblancAlinea
7. --- To quo interesas en la portuo pri qua\nl
me parolas, esas ke ol kontenas ne nur mili\cc
tala portuo, artifice aranjita per gigantala la\cc
bori, ma anke marveloza \VUgras{komercal portuo} si\cc
tuita en l'\VUgras{enflueyo} di granda fluvio.

%%K t10-08-1 p
\VUblancAlinea
8. --- La tera lango qua separas la du baseni\nl
esas fortifikita e defensata per serio de \VUgras{baterii}\nl
e \VUgras{bastioni} qui avancas sur la maro. On bezo\cc
nas specal permiso por promenar sur la \VUgras{rem}\cc
\VUgras{pari} \textsuperscript{(15)}. Me facile obtenabis ol, per mea on\cc
klulo, qua esas injenioro dil \VUgras{konstrukteyi} \textsuperscript{(18)}
\parplein
\end{VUpage}

% --- Feuillet 72 (folio 70) -----------------------------------------
\begin{VUpage}[72]{70}
%%K t10-08-1 p suite
\VUcontinue
e me iris vizitar la navi quin on konstruktas\nl
sub vasta hangari.

%%K t10-09-1 p
\VUblancAlinea
9. --- Mem me asistis la lanso di kuras-navo\nl
e me memoras la emoganta instanto quan la\nl
tota turbo sentis kande la grandega maso,\nl
abandonita a su ipsa, gliteskis sur sua bersi\cc
latra framo e venis flotacar sur l'aquo dil flu\cc
vio.

%%K t10-10-1 p
\VUblancAlinea
10. --- Por adirar la konstrukteyi, on trai\cc
ras la \VUgras{manovr-agro} \textsuperscript{(19)} ube la soldati di la gar\cc
nizono singladie venas exercar su, on pasas sur\nl
fera \VUgras{ponteto} \textsuperscript{(20)}, qua komunikigas la parad-es\cc
planado kun l'\VUgras{arsenalo} \textsuperscript{(21)}.

%%K t10-tit-4 sub
\VUsaut{3.0mm}
\VUpk{10.2pt}{40}{La Arsenalo e la doki.}
\VUsaut{3.0mm}

%%K t10-11-1 p
\VUblancAlinea
11. --- Kun mea onklulo, me vizitis ta granda\nl
konstrukturo plena de \VUgras{kanoni} \textsuperscript{(22)}, \VUgras{kugleqi} \textsuperscript{(23)}\nl
ed \VUgras{obusi}. Anke me vidis la \VUgras{metaliferio} \textsuperscript{(24)} en\nl
qua esas fabrikata la \VUgras{kaldieri} \textsuperscript{(25)} di la vapor\cc
navi.

%%K t10-12-1 p
\VUblancAlinea
12. --- Tre proxime, esas la \VUgras{doki} \textsuperscript{(26)} --- repar\cc
doki --- e la tre remarkinda \VUgras{flotacanta doki} \textsuperscript{(27)}\nl
en qui me povis vidar tre proxime, sur repa\cc
renda navo, la \VUgras{helico} \textsuperscript{(28)}, la \VUgras{kilio} \textsuperscript{(29)} e la \VUgras{gu}\cc
\VUgras{vernilo} \textsuperscript{(30)} di batelo.

%%K t10-13-1 p
\VUblancAlinea
13. --- La komercala portuo esas tote same\nl
studiinda kam la militala portuo, e me pasis\nl
multa hori gapante sur la kayi ube esas ama\cc
ragita la \VUgras{palet-navi} \textsuperscript{(31)} qui retrosequas la flu\cc
vio, e regardante funcionar la \VUgras{mastizo-kra}\cc
\VUgras{ni} \textsuperscript{(32)}, qui levas, quale plumin, enorma kar\cc
gaji.
\end{VUpage}

% --- Feuillet 73 (folio 71) -----------------------------------------
\begin{VUpage}[73]{71}
%%K t10-tit-5 sub
\VUsaut{4.0mm}
\VUpk{10.2pt}{40}{La Piloto.}
\VUsaut{3.0mm}

%%K t10-14-1 p
\VUblancAlinea
14. --- Me admiris la \VUgras{trans-Atlantiki} \textsuperscript{(34)} eki\cc
ranta la rado, kun lia \VUgras{flagi} \textsuperscript{(35)} en plena aero,\nl
duktata da habila \VUgras{piloto} \textsuperscript{(36)} qua marveloze sa\cc
vas alongirar la \VUgras{kanaleto} markizita per \VUgras{boyi} \textsuperscript{(40)}\nl
e \VUgras{balisi}, evitante la \VUgras{gabari} \textsuperscript{(41)} qui penigive di\cc
rektesas per lia unika seglo quadrata. Me di\cc
cernis sur l'\VUgras{avanajo} \textsuperscript{(37)} --- la pruo --- la ka\cc
\VUgras{bini} \textsuperscript{(39)} dil duesma klaso, e sur la \VUgras{dopajo} \textsuperscript{(38)}\nl
--- la pupo --- la saloni por la pasajeri dil unes\cc
ma klaso.

%%K t10-15-1 p
\VUblancAlinea
15. --- Kelkafoye me anke asistis la charjado\nl
di \VUgras{kargo-batelo} \textsuperscript{(42)} en qua jus esis amasigata\nl
la vari dil \VUgras{magazino} \textsuperscript{(43)} e di la \VUgras{mar-stacio}\cc
\VUgras{no} \textsuperscript{(44)}, adduktita per la multega treni dil \VUgras{kayo}\cc
\VUgras{lineo} \textsuperscript{(45)}.

%%K t10-tit-6 sub
\VUsaut{3.0mm}
\VUpk{10.2pt}{40}{La Remo-plezuro.}
\VUsaut{3.0mm}

%%K t10-16-1 p
\VUblancAlinea
16. --- Ofte me kanotagis en la \VUgras{flot-doko} \textsuperscript{(53)},\nl
en qua venas refujar la \VUgras{barki} \textsuperscript{(46)}, ed esis joyo\nl
por me manovrar la \VUgras{remilo} \textsuperscript{(47)}. Me acensis la\nl
\VUgras{ferdeko} \textsuperscript{(48)} di \VUgras{seglo-barko} \textsuperscript{(49)}, me klimis, per\nl
\VUgras{kordegi} \textsuperscript{(52)}, sur la \VUgras{masto} \textsuperscript{(51)} aden la \VUgras{yardi} \textsuperscript{(50)},\nl
pose me iteris mea promeno \VUgras{per yolo} \textsuperscript{(54)} inter\nl
pezoza \VUgras{pesko-barki} \textsuperscript{(55)}, ube la \VUgras{reti} \textsuperscript{(56)} sikes\cc
kas.

%%K t10-17-1 p
\VUblancAlinea
17. --- Sur la \VUgras{kayo} \textsuperscript{(58)} defilis koram me la\nl
maxim diversa \VUgras{vari} \textsuperscript{(59)}, kontenata en \VUgras{kesti} \textsuperscript{(60)}\nl
od en \VUgras{pakegi} \textsuperscript{(61)}. Li formacis la \VUgras{kargajo} \textsuperscript{(63)} dil\nl
\VUgras{barkasi}, e potenta \VUgras{krani} \textsuperscript{(62)} deprenis e depozis\nl
li sur \VUgras{kamioni} \textsuperscript{(64)} dum ke la \VUgras{kargo-veturisti}\nl
deklaris oli e pagis la taxi en la \VUgras{doganeyo} \textsuperscript{(65)}.
\end{VUpage}

% --- Feuillet 74 (folio 72) -----------------------------------------
\begin{VUpage}[74]{72}
%%K t10-18-1 p
\VUblancAlinea
18. --- Fine, kande me arivis a la \VUgras{turno}\cc
\VUgras{ponto} \textsuperscript{(66)} od a la \VUgras{sluzo} \textsuperscript{(69)}, pasante avan la\nl
\VUgras{drag-mashino} \textsuperscript{(68)} qua netigas la \VUgras{kanalo} \textsuperscript{(67)}, me\nl
desembarkis sur la warfo proxim la \VUgras{sema}\cc
\VUgras{foro} \textsuperscript{(70)}.

%%K t10-19-1 p
\VUblancAlinea
19. --- Sur l'altra latero di ta warfo ya venas\nl
abordar la \VUgras{shalupi} \textsuperscript{(71)} qui duktas a la tero l'ofi\cc
ciri dil eskadro. Esas admirinda spektaklo vi\cc
dar la maristi kadence levar sua remili, dum\nl
ke la barko portata dal \VUgras{ondo} \textsuperscript{(72)}, facile abordas\nl
la skalo dil desembarkeyo. Sur ica loko, on\nl
plu bone povas observar la \VUgras{mareo} e la movado\nl
dil \VUgras{fluxo} e \VUgras{refluxo} sur la \VUgras{plajo} \textsuperscript{(73)}.

%%K t10-tit-7 sub
\VUsaut{3.0mm}
\VUpk{10.2pt}{40}{La Klubeyo.}
\VUsaut{3.0mm}

%%K t10-20-1 p
\VUblancAlinea
20. --- Por retrovenar a la hemo, me uzis la\nl
\VUgras{elektrala tramvoyo} \textsuperscript{(74)} di qua la \VUgras{halteyo} \textsuperscript{(75)}\nl
esas proxim la \VUgras{fosto} pozita avan la klubo dil\nl
oficiri.

%%K t10-20-2 p
\VUblancAlinea
De sur la \VUgras{balkono} \textsuperscript{(76)} dil klubo, ornita per\nl
\VUgras{ankri} \textsuperscript{(77)}, kanoni e kuglegi, ofte me vidis asem\cc
blajo splendida de naval oficiri : \VUgras{lietnanti} \textsuperscript{(78)},\nl
kun sua espado, \VUgras{kapitani di artilrio} \textsuperscript{(79)} ed \VUgras{in}\cc
\VUgras{fantrio} \textsuperscript{(80)} kun sua \VUgras{sabro}. Dop li esis \VUgras{nava}\cc
\VUgras{no} \textsuperscript{(81)} qua adportis ad ili \VUgras{lorno}, kande li volis\nl
distingar to quo eventas fore.

%%K t10-21-1 p
\VUblancAlinea
21. --- Me anke vidis yuna \VUgras{sublietnanti} \textsuperscript{(82)}\nl
recevanta l'instrucioni di sua \VUgras{komandanto} \textsuperscript{(83)}\nl
o dil \VUgras{admiralo} \textsuperscript{(84)}, dum ke \VUgras{quatermastro} \textsuperscript{(85)}\nl
vartis por portar oli a la navo.

%%K t10-22-1 p
\VUblancAlinea
22. --- De ta balkono, la vidajo esas splen\cc
dida : on dominacas la tota plajo kun lua
\parplein
\end{VUpage}

% --- Feuillet 75 (folio 73) -----------------------------------------
\begin{VUpage}[75]{73}
%%K t10-22-1 p suite
\VUcontinue
\VUgras{tendi} \textsuperscript{(86)} e \VUgras{kabini} \textsuperscript{(87)}, ed on vidas, ye la horo\nl
dil balnado, la \VUgras{balneri} \textsuperscript{(88)} qui joyoze petulas\nl
avan la \VUgras{kazino} \textsuperscript{(89)}.

%%K t10-23-1 p
\VUblancAlinea
23. --- Ye l'extremajo di la plajo, la litoro\nl
formacas mikra \VUgras{peninsulo} \textsuperscript{(91)} \VUgras{rokoza}, ube la\nl
\VUgras{ondego} \textsuperscript{(92)} venas ruptar su e sur qua esas kons\cc
truktita \VUgras{faro} \textsuperscript{(93)}, nokte indikanta la voyo al\nl
naviganti. Ta peninsulo komunikas kun la\nl
tero nur per \VUgras{istmo} tre streta \textsuperscript{(90)}.

%%K t10-tit-8 sub
\VUsaut{3.0mm}
\VUpk{10.2pt}{40}{Generala vidajo di la litori.}
\VUsaut{3.0mm}

%%K t10-24-1 p
\VUblancAlinea
24. --- La \VUgras{klifo} \textsuperscript{(94)} extensas su, disfendita\nl
da la maro, qua kaprice desegnas en olu serio\nl
de \VUgras{bayi} \textsuperscript{(95)} e \VUgras{promontorii} (kapi), iganta lu\nl
maxim pitoreska.

%%K t10-24-2 p
\VUblancAlinea
Sur la \VUgras{platajo} \textsuperscript{(96)}, qua dominacas la \VUgras{(maro)}\cc
\VUgras{stretajo} \textsuperscript{(98)}, esas \VUgras{fuorto} \textsuperscript{(97)} tre importanta e\nl
kelka « villa »-i.

%%K t10-25-1 p
\VUblancAlinea
25. --- Ta stretajo separas del kontinento \VUgras{in}\cc
\VUgras{sulo} \textsuperscript{(99)} tre danjeroza, cirkondata da \VUgras{rifi} \textsuperscript{(100)} e\nl
\VUgras{brizanti}, ube barki e mem granda navi strandas\nl
ulfoye. Malgre la pronteso di la sokursi, e la\nl
rapida arivo dil \VUgras{salvo-kanoti}, ta \VUgras{naufraji} esas\nl
teroriganta, e, dum la \VUgras{tempesti} equinoxala,\nl
plura navi perisis kun homi e vari.

%%K t10-25-2 p
\VUblancAlinea
Malgre ta vicinajo, \VUgras{la portuo} esas tre fre\cc
quentata da la maristi.

\VUsaut{6.0mm}
\VUfilet{22mm}
\end{VUpage}
